\definecolor{usrbg}{RGB}{232,240,250}
\definecolor{usrfr}{RGB}{0,82,155}
\definecolor{thtbg}{RGB}{253,231,224}
\definecolor{thtfr}{RGB}{176,58,32}
\definecolor{txtfr}{RGB}{223,220,213}

\newcommand{\chipT}[1]{{\setlength{\fboxsep}{2.2pt}\setlength{\fboxrule}{.5pt}%
  \fcolorbox{thtfr}{thtbg}{\scriptsize\strut #1}}}
\newcommand{\chipX}[1]{{\setlength{\fboxsep}{2.2pt}\setlength{\fboxrule}{.5pt}%
  \fcolorbox{txtfr}{white}{\scriptsize\strut #1}}}

\begin{figure}[t]
\centering
\begin{tikzpicture}[
  turn/.style={rounded corners=2.5pt, draw=txtfr, line width=.6pt, fill=lstbg,
               inner xsep=4pt, inner ysep=4pt, font=\scriptsize},
  usr/.style={turn, fill=usrbg, draw=usrfr!55},
  arr/.style={-{Stealth[length=3.4pt,width=3pt]}, draw=black!42, line width=.6pt},
  ttl/.style={font=\scriptsize\bfseries, anchor=west},
  role/.style={font=\tiny, text=black!62, anchor=south},
  node distance=3.6mm,
]

\node[ttl] at (0,1.02) {Current Turn Injection};
\node[usr, anchor=west]  (c1) at (0,0) {user};
\node[turn, right=of c1] (c2) {\chipT{injected thought}};
\node[turn, right=of c2] (c3) {\chipX{text}};
\draw[arr] (c1) -- (c2);
\draw[arr] (c2) -- (c3);
\node[role] at (c2.north) {assistant};
\node[role] at (c3.north) {assistant continued};

\node[ttl] at ($(c3.east)+(1.5,1.02)$) {Past Turn Injection};
\node[usr, anchor=west]  (p1) at ($(c3.east)+(1.5,0)$) {user};
\node[turn, right=of p1] (p2) {\chipT{injected thought}\;\chipX{text}};
\node[usr,  right=of p2] (p3) {user};
\node[turn, right=of p3] (p4) {\chipX{text}};
\draw[arr] (p1) -- (p2);
\draw[arr] (p2) -- (p3);
\draw[arr] (p3) -- (p4);
\node[role] at (p2.north) {assistant};
\node[role] at (p4.north) {assistant};

\end{tikzpicture}
\looseness=-1 \caption{\textbf{Injection schemes exploiting in- and cross-session compatibility.} We find that
available form of thought injection depends on the provider and exact model. \textbf{Current-turn injection} places the thought in the current
assistant turn, so the model continues its visible answer straight from it; as of July 2026 it
is accepted by every GPT and Gemini model we tested and by the 4.5 generation of
Claude. In addition, Sonnet 4.5, Haiku 4.5, and all Gemini models accept
prefilling of the assistant turn's visible output. \textbf{Past-turn injection}
is available only against models that do not omit previous reasoning
blocks (e.g., Sonnet 5, Opus 4.8, Fable 5, and the GPT-5.6 series).}
\label{fig:injection_schemes}
\end{figure}
